\section{Discussion}\label{sec:discussion}

Three findings survive every caveat the first pass carries. The
negative trade shocks of the 2020s were regressive through prices, and
the positive ones offer no offsetting progressivity---partly because
they are small, partly because their gains arrive in categories
(clothing, footwear) whose budget shares are flat across the
distribution. The state's cushioning was margin-specific: automatic
where the shock hit earnings, discretionary where it hit prices---and
the discretionary instrument, capping a price, cushioned
proportionally rather than progressively, so the post-intervention
burden of the largest shock remained regressive. And the automatic
price-side stabiliser the welfare state does possess---indexation---is
calibrated to a world where prices move slowly: against a fast
terms-of-trade shock it delivered its help a year late, which for the
poorest households was the difference between the rulebook cushioning
the shock and amplifying it.

One asymmetry deserves to be stated as a finding about the policy
record rather than the method: Britain's negative trade shocks are
\emph{measured}---their first stages rest on published ex-post
estimates---while its positive trade shocks are \emph{promised},
resting on the government's own ex-ante projections. A comparative
household lens makes the asymmetry consequential: the measured
negative episodes are large and regressive; the projected positive
episodes are small and proportional. On the government's own numbers,
the distributional ledger of the 2020s trade policy portfolio is in
deficit, and nothing in the projected gains repays the measured
losses' incidence.

The limits are those declared throughout. First stages are imported,
with their identification and their controversies; the incidence is
first-order and static; the channels are not summed; the rulebook
decomposition labels legislated change as discretion without asking
what would have been legislated anyway. The first pass runs on decile
tables and representative households, not on household records---the
covariance between exposure and entitlement, the poverty and
inequality effects, and the population-weighted rulebook
decompositions all await the microsimulation stage, which the
verification work reported in Section~\ref{sec:roadmap} shows is
buildable with existing tools on existing data.

What the framework offers, finally, is a standing capability rather
than a one-off study: every future trade shock---a retaliation
schedule, the next agreement, the next terms-of-trade event---arrives
as one more declared first stage for the same second stage. A country
that intends to conduct an active trade policy while running a
means-tested welfare state needs exactly this join, and as of this
writing no institution, in Britain or elsewhere, maintains it.
